\documentclass[11pt]{article}
\usepackage[margin=1in]{geometry}
\usepackage{amsmath,amssymb,amsthm}
\usepackage{mathtools}
\usepackage{hyperref}

\newtheorem{theorem}{Theorem}
\newtheorem{lemma}[theorem]{Lemma}
\newtheorem{proposition}[theorem]{Proposition}
\newtheorem{corollary}[theorem]{Corollary}
\theoremstyle{definition}
\newtheorem{definition}[theorem]{Definition}
\newtheorem{example}[theorem]{Example}
\theoremstyle{remark}
\newtheorem*{remark}{Remark}

\newcommand{\R}{\mathbb{R}}
\newcommand{\C}{\mathbb{C}}
\newcommand{\F}{\mathbb{F}}
\newcommand{\id}{\mathrm{id}}
\newcommand{\tr}{\mathrm{tr}}
\DeclareMathOperator{\diag}{diag}
\DeclareMathOperator{\rank}{rank}

\title{Eigenvalues and Eigenvectors}
\author{Math Foundations --- Node 07}
\date{}

\begin{document}
\maketitle

\subsection*{First, an example}
Before any definitions, look at
\[
  A \;=\; \begin{pmatrix} 2 & 1 \\ 0 & 3 \end{pmatrix}.
\]
\emph{Read aloud:} ``A equals the upper-triangular two-by-two matrix with rows two one and zero three.''
Its characteristic polynomial is
\[
  \det(A - \lambda I) \;=\; (2-\lambda)(3-\lambda) \;=\; 0,
\]
\emph{Read aloud:} ``the determinant of A minus lambda I equals two-minus-lambda times three-minus-lambda equals zero.''
so the eigenvalues are $\lambda_1 = 2$ and $\lambda_2 = 3$. Solving
$(A - 2I)v = 0$ yields the eigenvector $v_1 = (1,0)^{\top}$, and
$(A - 3I)v = 0$ yields $v_2 = (1,1)^{\top}$. Each $v_i$ is mapped by $A$
to a pure rescaling: $A v_1 = 2 v_1$ and $A v_2 = 3 v_2$. The whole rest
of this lesson is the systematic theory behind that little calculation.

\section{Motivation}

A linear map $A : \R^n \to \R^n$ is, in general, a complicated object: it
mixes coordinates, rotates, shears, and rescales. The whole point of
eigentheory is to look for the directions on which $A$ acts trivially ---
not literally as the identity, but as pure scaling. Once such a direction
$v$ has been found, the action of $A$ along it is described by a single
number $\lambda$, the corresponding eigenvalue, via $A v = \lambda v$. If
enough such directions can be found to form a basis, then $A$ itself
factorises as a change of basis followed by a diagonal scaling. This is
the cleanest possible structural description of a linear map and it
underlies essentially every spectral technique in numerical analysis,
optimisation, and machine learning.

\section{Definitions}

Throughout, let $\F$ denote either $\R$ or $\C$ and let $A \in \F^{n\times n}$.

\begin{definition}[Eigenvalue, eigenvector]
A scalar $\lambda \in \F$ is an \emph{eigenvalue} of $A$ if there exists a
non-zero vector $v \in \F^{n}$ such that
\[
  A v \;=\; \lambda v.
\]
\emph{Read aloud:} ``A times v equals lambda times v.''
Any such $v$ is called an \emph{eigenvector} of $A$ associated with the
eigenvalue $\lambda$. The set of all eigenvectors for a fixed $\lambda$
together with the zero vector is a linear subspace of $\F^{n}$ called the
\emph{eigenspace} $E_\lambda = \ker(A - \lambda I)$. The set of all
eigenvalues of $A$ is the \emph{spectrum} $\sigma(A)$.
\end{definition}

\begin{remark}
The condition $v \ne 0$ is essential. Without it, every scalar $\lambda$
would satisfy $A \cdot 0 = \lambda \cdot 0$ trivially, and the definition
would carry no information.
\end{remark}

\begin{definition}[Characteristic polynomial]
The \emph{characteristic polynomial} of $A$ is
\[
  p_A(\lambda) \;=\; \det(A - \lambda I) \;\in\; \F[\lambda].
\]
\emph{Read aloud:} ``p sub A of lambda equals the determinant of A minus lambda I, an element of F adjoin lambda.''
It is a polynomial of degree exactly $n$ in $\lambda$ with leading term
$(-1)^{n} \lambda^{n}$.
\end{definition}

\begin{definition}[Algebraic and geometric multiplicity]
Let $\lambda_0 \in \sigma(A)$.
\begin{itemize}
  \item The \emph{algebraic multiplicity} $m_a(\lambda_0)$ is the
    multiplicity of $\lambda_0$ as a root of $p_A$.
  \item The \emph{geometric multiplicity} $m_g(\lambda_0)$ is the
    dimension of the eigenspace $E_{\lambda_0} = \ker(A - \lambda_0 I)$.
\end{itemize}
It is always the case that $1 \le m_g(\lambda_0) \le m_a(\lambda_0)$.
\end{definition}

\begin{definition}[Diagonalisability]
$A$ is \emph{diagonalisable} over $\F$ if there exist an invertible
$V \in \F^{n \times n}$ and a diagonal $\Lambda \in \F^{n\times n}$ such
that $A = V \Lambda V^{-1}$. Equivalently, $\F^{n}$ admits a basis of
eigenvectors of $A$, equivalently $m_a(\lambda) = m_g(\lambda)$ for every
$\lambda \in \sigma(A)$.
\end{definition}

\section{The fundamental theorem}

We now state and prove the main structural fact tying eigenvalues to
roots of the characteristic polynomial.

\begin{theorem}[Eigenvalues are roots of the characteristic polynomial]
\label{thm:char-poly}
Let $A \in \F^{n \times n}$. A scalar $\lambda \in \F$ is an eigenvalue
of $A$ if and only if
\[
  \det(A - \lambda I) \;=\; 0.
\]
\emph{Read aloud:} ``the determinant of A minus lambda I equals zero.''
\end{theorem}

\begin{proof}
We prove both implications.

\smallskip
\noindent\emph{($\Rightarrow$).} Suppose $\lambda$ is an eigenvalue of
$A$. By Definition~1 there exists a non-zero vector $v \in \F^{n}$ with
$A v = \lambda v$. Rearranging,
\[
  (A - \lambda I) v \;=\; A v - \lambda v \;=\; 0.
\]
\emph{Read aloud:} ``open paren A minus lambda I close paren times v equals A v minus lambda v equals zero.''
Thus the homogeneous linear system $(A - \lambda I) x = 0$ has the
non-trivial solution $x = v \ne 0$. A square matrix admits a non-trivial
kernel vector if and only if it is singular, equivalently iff its
determinant vanishes. Therefore $\det(A - \lambda I) = 0$.

\smallskip
\noindent\emph{($\Leftarrow$).} Conversely, suppose $\det(A - \lambda I)
= 0$. Then $A - \lambda I$ is a singular $n\times n$ matrix over $\F$,
so its kernel $\ker(A - \lambda I)$ is a non-trivial subspace of
$\F^{n}$ (it has positive dimension, by the rank--nullity theorem applied
to the linear map $x \mapsto (A - \lambda I) x$, whose rank is at most
$n - 1$). Choose any non-zero $v \in \ker(A - \lambda I)$. Then
$(A - \lambda I) v = 0$, i.e.\ $A v = \lambda v$, with $v \ne 0$. By
definition, $\lambda$ is an eigenvalue of $A$ with eigenvector $v$.
\end{proof}

\begin{corollary}
A matrix $A \in \C^{n\times n}$ has at least one eigenvalue in $\C$. In
fact it has exactly $n$ eigenvalues counted with algebraic multiplicity,
because $p_A \in \C[\lambda]$ splits completely over $\C$ by the
fundamental theorem of algebra.
\end{corollary}

\section{The spectral theorem (symmetric case)}

\begin{theorem}[Real spectral theorem]
\label{thm:spectral}
If $A \in \R^{n\times n}$ is symmetric ($A^{\top} = A$), then all
eigenvalues of $A$ are real, and there exists an orthogonal matrix
$Q \in \R^{n\times n}$ (i.e.\ $Q^{\top} Q = I$) and a diagonal matrix
$\Lambda = \diag(\lambda_1, \dots, \lambda_n)$ such that
\[
  A \;=\; Q \Lambda Q^{\top}.
\]
\emph{Read aloud:} ``A equals Q big-lambda Q transpose.''
\end{theorem}

\noindent We do not prove this here but record it because it is the
workhorse fact for covariance matrices, Gram matrices, and Laplacians.

\section{A small worked example}

\begin{example}
Let
\[
  A \;=\; \begin{pmatrix} 2 & 1 \\ 1 & 2 \end{pmatrix}.
\]
\emph{Read aloud:} ``A equals the two-by-two matrix with rows two one and one two.''
Then
\[
  p_A(\lambda) \;=\; \det \begin{pmatrix} 2-\lambda & 1 \\ 1 & 2-\lambda \end{pmatrix}
  \;=\; (2-\lambda)^{2} - 1 \;=\; \lambda^{2} - 4\lambda + 3
  \;=\; (\lambda-1)(\lambda-3).
\]
\emph{Read aloud:} ``p sub A of lambda equals the determinant of the matrix with rows two-minus-lambda one and one two-minus-lambda, which equals two-minus-lambda squared minus one, equals lambda squared minus four lambda plus three, equals lambda-minus-one times lambda-minus-three.''
Thus $\sigma(A) = \{1, 3\}$. Solving $(A - I)v = 0$ gives $v_1 = (1, -1)^{\top}/\sqrt{2}$, while $(A - 3I)v = 0$ gives $v_2 = (1, 1)^{\top}/\sqrt{2}$. These are orthonormal, as guaranteed by Theorem~\ref{thm:spectral}, and
\[
  A \;=\; \frac{1}{2} \begin{pmatrix} 1 & 1 \\ -1 & 1 \end{pmatrix}
  \begin{pmatrix} 1 & 0 \\ 0 & 3 \end{pmatrix}
  \begin{pmatrix} 1 & -1 \\ 1 & 1 \end{pmatrix}.
\]
\emph{Read aloud:} ``A equals one-half times the matrix with rows one one and minus-one one, times the diagonal matrix one-three, times the matrix with rows one minus-one and one one.''
\end{example}

\section{Where to go next}

The next step in this map is the singular value decomposition (SVD),
which extends eigendecomposition to rectangular and non-normal matrices
and underlies LoRA-style low-rank adaptation, PCA, and dimension
reduction in general.

\end{document}
